\chapter{Creatinine Clearance}

\section*{What Is This Test?}
Creatinine clearance estimates the kidney's ability to remove creatinine from the blood. It combines a timed urine collection with a blood creatinine measurement and uses urine volume and creatinine concentration to estimate clearance, reported in mL/min.

\section*{Alternative Names}
Measured Creatinine Clearance; 24-Hour Creatinine Clearance; Renal Clearance Test.

\section*{Why This Test Is Ordered}
It may be used when a measured estimate of kidney filtration is needed, when creatinine-based estimated GFR is unreliable, or when drug dosing and kidney function require additional assessment. Modern eGFR equations are usually preferred for routine chronic kidney disease evaluation.

\section*{How This Test Is Performed}
The patient empties the bladder at the start time and discards that urine, then collects every urine specimen for the timed interval, commonly 24 hours. A blood sample is collected during the collection period. The laboratory calculates clearance from urine creatinine, plasma creatinine, and urine flow rate, with body-surface-area adjustment when appropriate.

\section*{Preparation, Risks, and Limitations}
Follow the laboratory's collection instructions and keep the container refrigerated if instructed. Do not discard samples or add urine from outside the timed period. There is no procedural risk beyond venipuncture, but incomplete collection can make the result falsely low. Muscle mass, diet, medications, unstable kidney function, and high urine output also affect accuracy.

\section*{Understanding Results}
Lower clearance generally indicates reduced filtration, but the result must be interpreted with age, body size, muscle mass, hydration, and clinical history. A result that conflicts with the eGFR or clinical picture should be reviewed rather than used in isolation.

\section*{References}
National Kidney Foundation guidance on GFR estimation and measured kidney function.

\clearpage
\section*{Understanding Results and Follow-up}
The laboratory report should identify the specimen, method, units, reference interval or decision limit, and whether the result is positive, negative, abnormal, or inconclusive. Reference intervals vary with age, sex, pregnancy, specimen type, assay, and laboratory; a result outside a range is not automatically a diagnosis, and a result within a range does not always exclude disease. Discuss unexpected results with the ordering clinician, who can decide whether repeat or confirmatory testing, treatment, or referral is appropriate. Do not change medicines or delay urgent care based on an isolated result.


\section*{Regulatory Status and Authorization Date}
Regulatory status applies to the specific assay, instrument, manufacturer, and intended use rather than to the test name alone. No single approval date can be assigned to this test category. For a commercial assay, verify the current product in the relevant regulator's database: the U.S. Food and Drug Administration (FDA) clearance, approval, or authorization; India's Central Drugs Standard Control Organisation (CDSCO) licence or permission; the European Union In Vitro Diagnostic Regulation (EU IVDR) CE certificate issued through the applicable conformity-assessment route; or the United Kingdom Medicines and Healthcare products Regulatory Agency (MHRA) registration and UKCA/CE status. Laboratory-developed tests may not have an individual FDA, CDSCO, EU IVDR, or MHRA approval date.
\clearpage
